\chapter{Bloco 7 -- \emph{Bin packing}}
\label{cap:bloco7}

Empacotar \textbf{todos} os itens no \textbf{menor número de caixas} de capacidade $b$. É a mochila ``ao
contrário'': agora todo item entra, e o que queremos minimizar é o número de caixas.

\paragraph{Os dados.} Dez itens de tamanhos $w=(60,36,53,51,58,48,47,48,53,35)$, caixas de capacidade $b=100$. A soma
dos tamanhos é 489, então o \textbf{limite de volume} é $\lceil489/100\rceil=5$ caixas.

\begin{discussao}
\begin{enumerate}
  \item \emph{Quais variáveis?} $x_{ij}=1$ se o item $j$ vai para a caixa $i$; $y_i=1$ se a caixa $i$ é usada.
  \item \emph{Quantas caixas considerar?} No pior caso, uma por item: $n$ caixas ($i=1,\dots,n$).
  \item \emph{As duas famílias de restrições?} Todo item vai para \textbf{exatamente uma} caixa, e a capacidade é
  respeitada, \textbf{mas só se a caixa for usada}. ``Só se usada'' é o Bloco 1 de novo:
  $\sum_jw_jx_{ij}\le b\,y_i$.
  \item \emph{Dá para usar 5 caixas?} Não; veja o argumento abaixo.
\end{enumerate}
\end{discussao}

\section{O modelo}
\[
\begin{aligned}
\min\;& \sum_{i=1}^{n} y_i\\
\sa & \sum_{i=1}^{n} x_{ij} = 1 &&\forall j &&\text{(todo item em exatamente uma caixa)}\\
& \sum_{j=1}^{n} w_j x_{ij} \le b\,y_i &&\forall i &&\text{(capacidade, só se a caixa é usada)}\\
& x_{ij},\,y_i\in\B.
\end{aligned}
\]
Repare que $b\,y_i$ faz dois papéis ao mesmo tempo: se $y_i=1$, limita a ocupação a $b$; se $y_i=0$, proíbe
colocar qualquer item na caixa $i$ (é uma restrição de ativação com $M=b$).

\begin{codigo}
def modelo_binpacking(w, b, simetria=False):
    n = len(w)
    prob = pulp.LpProblem("binpacking", pulp.LpMinimize)
    y = {i: pulp.LpVariable(f"y_{i}", cat="Binary") for i in range(n)}
    x = {(i, j): pulp.LpVariable(f"x_{i}_{j}", cat="Binary") for i in range(n) for j in range(n)}
    prob += pulp.lpSum(y.values()), "n_caixas"                                   # TODO 1
    for j in range(n):                                                           # TODO 2
        prob += pulp.lpSum(x[i, j] for i in range(n)) == 1, f"item_{j}_em_uma_caixa"
    for i in range(n):                                                           # TODO 3
        prob += pulp.lpSum(w[j]*x[i, j] for j in range(n)) <= b*y[i], f"capacidade_{i}"
    if simetria:                                                                 # desafio
        for i in range(n-1):
            prob += y[i] >= y[i+1], f"simetria_{i}"
    return prob, x, y
\end{codigo}

\section{Resultado: 6 caixas}

\begin{center}
\begin{tikzpicture}[x=1.45cm,y=0.032cm]
  \foreach \cx/\a/\b/\la/\lb in {1/60/35/60/35, 2/58/0/58/, 3/36/53/36/53, 4/51/48/51/48, 5/47/53/47/53, 6/48/0/48/}{
    \draw[thick] (\cx-0.35,0) rectangle (\cx+0.35,100);
    \fill[azul!35] (\cx-0.35,0) rectangle (\cx+0.35,\a);
    \node[font=\small] at (\cx,\a/2) {\la};
    \ifnum\b>0
      \fill[verde!35] (\cx-0.35,\a) rectangle (\cx+0.35,\a+\b);
      \node[font=\small] at (\cx,\a+\b/2) {\lb};
    \fi
    \node[below,font=\footnotesize] at (\cx,0) {caixa \cx};
  }
  \draw[dashed,red] (0.5,100) -- (6.5,100) node[right,font=\footnotesize] {$b=100$};
\end{tikzpicture}
\end{center}
Ocupações: 95, 58, 89, 99, 100 e 48. A heurística FFD (\emph{First-Fit Decreasing}) do notebook também usa 6 caixas,
e a relaxação linear vale $4{,}89$.

\paragraph{Por que é impossível usar 5 caixas?} Um argumento curto, só com aritmética:
\begin{enumerate}
  \item Há \textbf{cinco itens maiores que 50}: 60, 58, 53, 53 e 51. Dois deles nunca cabem juntos, então cada um
  ocupa uma caixa diferente. Com 5 caixas, cada caixa teria exatamente um item grande.
  \item As folgas dessas caixas seriam $40,\ 42,\ 47,\ 47$ e $49$.
  \item Os itens pequenos são 48, 48, 47, 36 e 35. Dois pequenos juntos somam pelo menos $35+36=71>49$, então cada
  caixa recebe \emph{no máximo um} pequeno.
  \item Os \textbf{dois itens de tamanho 48} precisam de folga $\ge48$, mas só a caixa do 51 tem folga 49. Um dos 48
  fica sem lugar.
\end{enumerate}
Logo são necessárias pelo menos 6 caixas, e o solver encontrou 6: \textbf{6 é ótimo}.

\begin{outros}[limitantes para o \emph{bin packing}]
O passo 1 do argumento é a ideia por trás dos limitantes inferiores clássicos de Martello e Toth: itens maiores que
metade da capacidade não podem dividir caixa entre si e precisam ser contados separadamente, o que melhora o limite
de volume $\lceil\sum w_j/b\rceil$~\cite{martellototh1990,delorme2016}. É um exemplo de como raciocínio
combinatório complementa a relaxação linear.
\end{outros}

\section{Reflexões: relaxação fraca e simetria}

\paragraph{Relaxação fraca.} Na relaxação, o solver pode usar $y_i=\sum_jw_jx_{ij}/b$, ``meia caixa'' aqui e ali, e
o valor fica exatamente $\sum_jw_j/b=4{,}89$. A relaxação não ``enxerga'' que itens grandes não se combinam.

\begin{outros}[a formulação ``de Kantorovich'' e alternativas]
O modelo com $x_{ij}$ e $y_i$ é atribuído a Kantorovich~\cite{kantorovich1960}. Os surveys de Valério de
Carvalho~\cite{carvalho2002} e de Delorme, Iori e Martello~\cite{delorme2016} registram que sua relaxação linear é
fraca (vale apenas o limite de volume) e que formulações baseadas em \emph{padrões de corte}, como a de Gilmore e
Gomory~\cite{gilmoregomory1961}, têm relaxações muito mais fortes, ao custo de um número exponencial de variáveis.
\end{outros}

\paragraph{Simetria.} As caixas são \textbf{idênticas}: trocar os rótulos de duas caixas gera uma solução
``diferente'' para o solver, mas igual na prática. Uma solução com $k$ caixas tem $k!$ cópias equivalentes, e o
\emph{branch-and-bound} pode perder tempo explorando todas.

\paragraph{Quebra de simetria.} Exigir que as caixas sejam usadas \emph{em ordem}:
\[
y_i\ge y_{i+1}\qquad(i=1,\dots,n-1).
\]
Isso não muda o ótimo (sempre dá para renumerar as caixas usadas como $1,2,\dots,k$), mas elimina as cópias. No
experimento do notebook com 14 itens, sem a quebra de simetria o CBC bateu o limite de 30 s; com ela, resolveu em
0,15 s (ambos com 6 caixas). Os tempos variam com máquina e solver, mas a tendência é clara.

\begin{atencao}[erros típicos do Bloco 7]
\begin{itemize}
  \item Escrever a capacidade como $\sum_jw_jx_{ij}\le b$ (sem $y_i$): nada liga $y$ aos itens, e o solver zera
  todos os $y_i$, devolvendo ``0 caixas''.
  \item Usar $\sum_ix_{ij}\le1$ no lugar de $=1$: o modelo ``descarta'' itens e usa 0 caixas.
  \item Com \texttt{timeLimit}, o PuLP pode reportar ``Optimal'' para uma solução que é apenas viável. Desconfie de
  tempos que coincidem com o limite.
\end{itemize}
\end{atencao}
